\section{Triton Softmax：行归约操作}

\subsection{从逐元素到行归约}

前面的 GeLU 是\textbf{逐元素}操作——每个输出只依赖一个输入。Softmax 则是\textbf{行归约}操作——每个输出依赖\textbf{整行}的所有输入：

$$\text{softmax}(x_i) = \frac{e^{x_i - \max(\mathbf{x})}}{\sum_j e^{x_j - \max(\mathbf{x})}}$$

\begin{itemize}
    \item $x_i$：第 $i$ 个元素
    \item $\max(\mathbf{x})$：整行的最大值（数值稳定性）
    \item $\sum_j$：整行的求和
\end{itemize}

这意味着需要跨多个元素\textbf{通信}（求最大值、求和），不能简单地让每个线程独立工作。

\subsection{朴素 Softmax 实现与内存分析}

\begin{lstlisting}[caption={朴素 Softmax 实现及内存读写分析}]
def manual_softmax(x):
    M, N = x.shape
    x_max = x.max(dim=1)[0]         # MN reads, M writes
    x = x - x_max[:, None]          # MN+M reads, MN writes
    numerator = torch.exp(x)        # MN reads, MN writes
    denominator = numerator.sum(dim=1)  # MN reads, M writes
    y = numerator / denominator[:, None]  # MN+M reads, MN writes
    # Total: 5MN + 2M reads, 3MN + 2M writes
    return y
\end{lstlisting}

\begin{importantbox}{朴素 Softmax 的内存开销}
朴素实现需要 $5MN + 2M$ 次读和 $3MN + 2M$ 次写（全局内存访问）。
而理论最优只需 $MN$ 次读和 $MN$ 次写——\textbf{潜在的 4 倍加速}！
\end{importantbox}

\subsection{Triton Softmax 实现}

核心设计思想：\textbf{每个 Block 处理一行}。如果行的宽度能放入一个 Block（即放入 SM 的 Shared Memory），那么整行的 max、sum、归一化都可以在 Block 内完成，只需一次全局内存读和一次写。

\subsubsection{包装器函数}

\begin{lstlisting}[caption={Triton Softmax 包装器}]
def triton_softmax(x):
    y = torch.empty_like(x)
    M, N = x.shape
    block_size = triton.next_power_of_2(N)  # 列数向上取整到 2 的幂
    num_blocks = M                           # 每行一个 Block

    triton_softmax_kernel[(M,)](
        x_ptr=x, y_ptr=y,
        x_row_stride=x.stride(0),
        y_row_stride=y.stride(0),
        num_cols=N, BLOCK_SIZE=block_size
    )
    return y
\end{lstlisting}

\subsubsection{内核函数}

\begin{lstlisting}[caption={Triton Softmax 内核——每个 Block 处理一行}]
@triton.jit
def triton_softmax_kernel(x_ptr, y_ptr,
    x_row_stride, y_row_stride,
    num_cols, BLOCK_SIZE: tl.constexpr):

    # 当前处理哪一行
    row_idx = tl.program_id(0)
    col_offsets = tl.arange(0, BLOCK_SIZE)

    # 读取整行到 SM 内部
    x_start_ptr = x_ptr + row_idx * x_row_stride
    x_ptrs = x_start_ptr + col_offsets
    x_row = tl.load(x_ptrs,
        mask=col_offsets < num_cols,
        other=float("-inf"))      # 超出列数的部分用 -inf 填充

    # 在 SM 内部完成所有计算
    x_row = x_row - tl.max(x_row, axis=0)  # 数值稳定
    numerator = tl.exp(x_row)
    denominator = tl.sum(numerator, axis=0)
    y_row = numerator / denominator

    # 写回全局内存
    y_start_ptr = y_ptr + row_idx * y_row_stride
    y_ptrs = y_start_ptr + col_offsets
    tl.store(y_ptrs, y_row, mask=col_offsets < num_cols)
\end{lstlisting}

\begin{importantbox}{Triton Softmax 的设计精髓}
\begin{enumerate}
    \item \textbf{一行 = 一个 Block}：利用 Block 内的 Shared Memory 完成所有归约操作
    \item \textbf{只需一次全局读写}：整行数据一次性加载到 SM，计算完成后一次性写回
    \item \textbf{超出列数的元素用 $-\infty$ 填充}：确保 \texttt{tl.max} 和 \texttt{tl.exp} 的正确性（$e^{-\infty} = 0$）
\end{enumerate}
\end{importantbox}

\begin{knowledgebox}{Block Size 为什么是 next\_power\_of\_2(N)？}
Triton 要求 Block Size 必须是 2 的幂（这是硬件要求）。使用 \texttt{triton.next\_power\_of\_2(N)} 确保 Block 足够大以容纳整行，多余的位置用 mask 和 \texttt{other=float("-inf")} 处理。
\end{knowledgebox}

\subsection{Softmax 的读写次数详细分析}

对 $M \times N$ 矩阵的逐行 Softmax，我们可以精确分析全局内存的读写次数：

\begin{table}[H]
\centering
\begin{tabular}{lcc}
\toprule
\textbf{操作} & \textbf{读次数} & \textbf{写次数} \\
\midrule
\texttt{x.max(dim=1)} & $MN$ & $M$ \\
\texttt{x - x\_max[:, None]} & $MN + M$ & $MN$ \\
\texttt{torch.exp(x)} & $MN$ & $MN$ \\
\texttt{numerator.sum(dim=1)} & $MN$ & $M$ \\
\texttt{numerator / denom[:, None]} & $MN + M$ & $MN$ \\
\midrule
\textbf{合计} & $5MN + 2M$ & $3MN + 2M$ \\
\textbf{融合最优} & $MN$ & $MN$ \\
\textbf{加速比} & \multicolumn{2}{c}{$\sim 4\times$} \\
\bottomrule
\end{tabular}
\caption{朴素 vs 融合 Softmax 的全局内存访问分析}
\end{table}

这个分析解释了为什么融合版本的 Softmax 加速比约为 2--3 倍（理论上限约 4 倍）——实际中融合内核还有一些不可避免的开销（如 Block 调度、掩码处理）。

\subsection{Softmax 完整性能对比}

\begin{table}[H]
\centering
\begin{tabular}{lcc}
\toprule
\textbf{实现} & \textbf{耗时 (ms)} & \textbf{相对于 PyTorch} \\
\midrule
Manual（朴素 PyTorch） & $\sim$3.7 & 2.5x 慢 \\
Triton（自写内核） & $\sim$1.9 & 1.3x 慢 \\
PyTorch（内置） & $\sim$1.5 & 1.0x \\
torch.compile & $\sim$1.3 & 0.87x（更快！）\\
\bottomrule
\end{tabular}
\caption{Softmax 四种实现的性能对比}
\end{table}

值得注意的是：
\begin{itemize}
    \item \texttt{torch.compile} 在 Softmax 上甚至超越了 PyTorch 内置实现——因为它知道输入的 shape，可以做更针对性的优化
    \item 朴素实现的 3.7ms 到融合实现的 1.3--1.9ms，约 2--3x 加速，不如 GeLU 的 8x——因为 Softmax 本身就有归约操作，理论上限约 4x
\end{itemize}

\begin{warningbox}{本实现假设行宽 $\leq$ Block Size}
当矩阵的列数非常大（超过 SM 的 Shared Memory 容量）时，单个 Block 无法容纳整行。此时需要更复杂的多 Block 协作方案，如 FlashAttention 中使用的 online softmax 和 tiling 策略。
\end{warningbox}

\subsection{本章小结}

Softmax 展示了 Triton 处理行归约操作的能力。核心设计是``一行 = 一个 Block''：整行数据加载到 SM 内部，在高速内存中完成所有归约计算，最后一次性写回。当行宽适中时，这种设计非常高效。

%% ========== Section 10: 矩阵乘法 ==========

